% Part: first-order-logic
% Chapter: beyond
% Section: many-sorted-logic

\documentclass[../../../include/open-logic-section]{subfiles}

\begin{document}

\olfileid{fol}{byd}{msl}

\olsection{Logica met verschillende soorten objecten}

In de eerste-ordelogica lopen alle variabelen en kwantoren over één
!!{domain}. Soms wil je objecten echter verdelen over meerdere
!!{domain}s die niet door elkaar lopen. Je kunt dan bijvoorbeeld één
!!a{domain} voor getallen nemen, één !!a{domain} voor meetkundige
objecten, één !!a{domain} voor functies van getallen naar getallen en
één !!a{domain} voor abelse groepen, enzovoort.

Meersoortige logica is het formele systeem voor zo'n indeling. Eerst
maak je een lijst met ``soorten''. De soort van een object zegt in welk
``!!{domain}'' dat object thuishoort. Voor elke soort gebruik je eigen
!!{variable}s en kwantoren, en meestal ook een eigen !!{identity}, dus
een eigen gelijkheidsrelatie. Bij elke functie en relatie leg je
bovendien vast welke soorten objecten op elke argumentplaats mogen
staan: de functie of relatie is dan naar soort ``getypeerd''. Voor het
overige blijven de gewone regels van de eerste-ordelogica gelden. De
regels voor kwantoren worden voor elke soort apart gegeven.

Neem als voorbeeld onderzoek naar internationale betrekkingen. We
kiezen twee soorten objecten: Franse burgers en Duitse burgers. Voor
de Franse burgers nemen we een unaire relatie ``drinkt wijn''. Dat is
een relatie die van één object iets zegt. Voor de Duitse burgers nemen
we op dezelfde manier de unaire relatie ``eet worst''. Daarnaast nemen
we de binaire relatie ``vormen samen een echtpaar met verschillende
nationaliteiten''. Die relatie verbindt twee objecten: het eerste
argument moet een Franse burger zijn en het tweede een Duitse burger.
Laat de variabelen $a$, $b$, $c$ over Franse burgers lopen en $x$, $y$, $z$
over Duitse burgers. Dan zegt de formule
\[
\lforall[a][\lforall x][(\Atom{\Obj{MarriedTo}}{a,x} \lif
(\Atom{\Obj{DrinksWine}}{a} \lor \lnot \Atom{\Obj{EatsWurst}}{x})]]
\]
het volgende: voor iedere Franse burger en iedere Duitse burger geldt
dat, als zij met elkaar getrouwd zijn, de Franse burger wijn drinkt of
de Duitse burger geen worst eet.

Je kunt deze logica op een natuurlijke manier omzetten in gewone
eerste-ordelogica. Eerst stop je alle objecten uit de verschillende
!!{domain}s in één groot eerste-orde-!!{domain}. Daarna voeg je voor
iedere soort een unair !!{predicate} toe dat aangeeft welke objecten
tot die soort behoren. Tot slot zorg je ervoor dat een kwantor alleen
over objecten van de bedoelde soort loopt. De eerste-ordetaal voor ons
voorbeeld krijgt zo, naast de relaties die we al noemden, de unaire
!!{predicate}s ``$\Obj{German}$'' en ``$\Obj{French}$''. Bij de andere
relaties laten we de soortvereisten zelf weg. Neem nu een
soortgebonden kwantor $\lforall[x][!A]$, waarbij $x$ !!a{variable} van de
Duitse soort is. Die vertalen we als
\[
\lforall[x][(\Atom{\Obj{German}}{x} \lif !A)].
\]
We moeten ook axioma's toevoegen die vastleggen dat de soorten uit
elkaar blijven, bijvoorbeeld
$\lforall[x][\lnot (\Atom{\Obj{German}}{x} \land
  \Atom{\Obj{French}}{x})]$. Verder zijn er axioma's nodig die ervoor
zorgen dat ``drinkt wijn'' alleen geldt voor objecten die aan het
predicaat $\Atom{\Obj{French}}{x}$ voldoen, enzovoort. Met deze afspraken
en axioma's kun je aantonen dat meersoortige !!{sentence}s in
eerste-orde-!!{sentence}s kunnen worden vertaald. Hetzelfde geldt voor
meersoortige !!{derivation}s en eerste-orde-!!{derivation}s.
Meersoortige !!{structure}s kun je eveneens in de bijbehorende
eerste-orde-!!{structure}s ``vertalen'', en ook weer terug. Daardoor
geldt voor de meersoortige logica eveneens een volledigheidsstelling.

\end{document}

